% SPDX-FileCopyrightText: 2026 kimihiro64
% SPDX-License-Identifier: Apache-2.0 OR CC-BY-4.0

\documentclass[11pt]{amsart}

\usepackage[T1]{fontenc}
\usepackage{lmodern}
\usepackage[nopatch=footnote]{microtype}
\usepackage[margin=1.1in]{geometry}
\usepackage{amsmath,amssymb,amsthm,mathtools}
\usepackage{booktabs}
\usepackage{xcolor}
\usepackage{hyperref}

\definecolor{linkblue}{HTML}{174A7E}
\hypersetup{
  colorlinks=true,
  linkcolor=linkblue,
  citecolor=linkblue,
  urlcolor=linkblue,
  pdftitle={Erdos 690: Prime-Factor Density Unimodality},
  pdfauthor={kimihiro64}
}

\newtheorem{theorem}{Theorem}[section]
\newtheorem{lemma}[theorem]{Lemma}
\newtheorem{proposition}[theorem]{Proposition}
\newtheorem{corollary}[theorem]{Corollary}
\theoremstyle{definition}
\newtheorem{definition}[theorem]{Definition}
\theoremstyle{remark}
\newtheorem{remark}[theorem]{Remark}

\title{Erdos 690: Prime-Factor Density Unimodality}
\author{kimihiro64}
\date{\today}

\begin{document}

\begin{abstract}
Lean formalization of the finite density recursion and non-unimodality certificates for ordered distinct prime factors.
\end{abstract}

\maketitle

\begin{center}
\small
Copyright \textcopyright\ 2026 kimihiro64. This paper is available under
the \href{https://www.apache.org/licenses/LICENSE-2.0}{Apache License 2.0} or
\href{https://creativecommons.org/licenses/by/4.0/}{CC BY 4.0}, at your option.
\end{center}

\section{Introduction}

Explain the problem's significance, research audience, source relationship,
and the exact scope of the result. Do not claim the project is complete while
the public Challenge/Solution surface is still the template sentinel.

\section{Statement of the problem}

\textbf{Unimodality of the density of the kth smallest distinct prime factor.} For each integer k at least 1 and prime p, let d\_k(p) be the natural density of positive integers whose kth smallest distinct prime divisor is p. Classify those k for which the prime-indexed sequence p mapsto d\_k(p) is unimodal.

Define every nonstandard object and state all quantifiers, hypotheses,
endpoints, and exceptional cases.

\section{Relation to the literature}

Give a precise source and prior-formalization account. Distinguish direct
formalization, adaptation, independent proof, background, and original work.

\section{Proof architecture}

Describe the mathematical reductions and how they compose. Replace this
scaffold with the actual proof, not a development chronology.

\section{Formalization and verification}

Map each headline theorem to its Lean declaration. Describe dependency pins,
the Challenge/Solution trust boundary, kernel checks, Comparator/NanoDa replay,
and any exact finite certificates or computations.

\section{Limitations and review status}

State known divergences, omitted material, automation use, and the review that
actually occurred. Mechanical checks are not independent expert review.

\section*{Acknowledgements}

Credit mathematical sources, libraries, tools, contributors, and reviewers.

\begin{thebibliography}{9}

\bibitem{PrimarySource}
Replace with the complete primary-source citation.

\end{thebibliography}

\end{document}
